% Use Case Diagram aplikasi NADI
% Kompilasi: pdflatex use-case.tex
\documentclass[border=12pt]{standalone}
\usepackage[T1]{fontenc}
\usepackage{lmodern}
\usepackage{tikz}
\usetikzlibrary{shapes.geometric,arrows.meta,positioning}

\renewcommand{\familydefault}{\sfdefault}

\begin{document}
\begin{tikzpicture}[
    uc/.style={draw, thick, ellipse, align=center, font=\small,
               minimum width=4.4cm, minimum height=1.25cm, inner sep=1pt, fill=white},
    garis/.style={thick},
    relasi/.style={thick, dashed, -{Straight Barb[length=6pt]}},
    stereotip/.style={font=\footnotesize, fill=white, inner sep=1pt}
  ]

  % Aktor berupa orang-orangan
  \newcommand{\aktor}[4]{%
    \begin{scope}[shift={(#2)}, thick]
      \draw (0,0.75) circle (0.25);
      \draw (0,0.5) -- (0,-0.3);
      \draw (-0.45,0.25) -- (0.45,0.25);
      \draw (0,-0.3) -- (-0.35,-0.9);
      \draw (0,-0.3) -- (0.35,-0.9);
      \node[align=center, font=\small, below] at (0,-0.95) {#3};
      \coordinate (#1) at (#4);
    \end{scope}}

  % Batas sistem
  \draw[thick] (2.6,0.9) rectangle (15.0,16.5);
  \node[font=\bfseries, anchor=north] at (8.8,16.4) {Aplikasi NADI};

  % Use case utama (kolom kiri)
  \node[uc] (pilih)   at (5.4,15.0) {Memilih wilayah};
  \node[uc] (data)    at (5.4,12.8) {Menampilkan data\\iklim wilayah};
  \node[uc] (insight) at (5.4,10.6) {Melihat tren historis\\(Climate Insight)};
  \node[uc] (banding) at (5.4,8.6)  {Membandingkan\\dua wilayah};
  \node[uc] (chat)    at (5.4,6.6)  {Bertanya di\\Chat AI};
  \node[uc] (aksi)    at (5.4,4.4)  {Melihat rekomendasi\\aksi iklim};
  \node[uc] (atur)    at (5.4,2.2)  {Mengatur koneksi\\dan API key};

  % Use case pendukung (kolom kanan)
  \node[uc] (unduh)   at (12.0,12.8) {Mengunduh data\\Open-Meteo};
  \node[uc] (jelas)   at (12.0,10.6) {Membaca penjelasan AI};
  \node[uc] (hapus)   at (12.0,5.0)  {Menghapus\\riwayat chat};
  \node[uc] (uji)     at (12.0,2.2)  {Menguji koneksi\\database / Gemini};

  % Aktor
  \aktor{pengguna}{0,8.6}{Pengguna}{0.5,0.1};
  \aktor{openmeteo}{17.6,12.8}{Open-Meteo\\API}{-0.5,0.1};
  \aktor{gemini}{17.6,8.0}{Google\\Gemini API}{-0.5,0.1};

  % Asosiasi pengguna
  \foreach \u in {pilih,data,insight,banding,chat,aksi,atur}
    \draw[garis] (pengguna) -- (\u.west);

  % Asosiasi sistem eksternal
  \draw[garis] (unduh.east) -- (openmeteo);
  \draw[garis] (jelas.east) -- (gemini);
  \draw[garis] (chat) -- (gemini);

  % include / extend
  \draw[relasi] (data) -- node[stereotip, above=2pt] {\guillemotleft include\guillemotright} (unduh);
  \draw[relasi] (insight) -- node[stereotip, above=2pt] {\guillemotleft include\guillemotright} (jelas);
  \draw[relasi] (hapus) -- node[stereotip, sloped, below=2pt] {\guillemotleft extend\guillemotright} (chat);
  \draw[relasi] (uji) -- node[stereotip, above=2pt] {\guillemotleft extend\guillemotright} (atur);

\end{tikzpicture}
\end{document}
